% Adapted from https://composingprograms.com/3ed/examples/#link-py
For the following questions, use this dataclass-based linked list.
A \lstinline{LinkedList[T]} is either a \lstinline{Link[T]} or the empty tuple
\lstinline{()}. The type variable \lstinline{T} describes the element type.

\begin{lstlisting}
from __future__ import annotations
from dataclasses import dataclass

type LinkedList[T] = Link[T] | tuple[()]

@dataclass
class Link[T]:
    first: T
    rest: LinkedList[T] = ()

    def __str__(self) -> str:
        return format_link(self)

def format_link[T](s: Link[T]) -> str:
    string = '(' + str(s.first)
    remaining = s.rest
    while isinstance(remaining, Link):
        string += ' ' + str(remaining.first)
        remaining = remaining.rest
    assert remaining == (), 'Not a valid linked list'
    return string + ')'
\end{lstlisting}

Construct a list with \lstinline{Link(1, Link(2))}. Read its elements with
\lstinline{first} and \lstinline{rest}. At the end, \lstinline{rest} is
\lstinline{()}, which has no \lstinline{first} or \lstinline{rest} attributes.
\lstinline{print} uses parentheses around the elements; entering a \lstinline{Link}
at the prompt shows the dataclass representation with field names.

\begin{meta}
\textbf{Teaching the interface:} Establish construction, field access, and the
empty case first. Treat \lstinline{format_link} as supplied display behavior;
students need not implement it. In the questions, reason about element values
and shorter lists. If using diagrams, draw the sequence structure without
asking students to distinguish shared nodes from equal-valued nodes.

This definition uses Python 3.12 or later. The future import also supports
these annotations on Python 3.12--3.13. Reference:
\href{https://composingprograms.com/3ed/compound-data/}{Composing Programs, Compound Data}.
\end{meta}
